\documentclass[12pt,a4paper]{article}

\usepackage[margin=2.5cm]{geometry}
\usepackage{booktabs}
\usepackage{amsmath}
\usepackage{amssymb}
\usepackage[hidelinks]{hyperref}
\usepackage{caption}
\usepackage{array}
\usepackage{microtype}
\usepackage{parskip}
\usepackage{setspace}
\usepackage[T1]{fontenc}
\usepackage{lmodern}
\usepackage{xcolor}
\usepackage{tabularx}

\setstretch{1.15}

\title{\textbf{Credit Card Default Prediction}\\[0.4em]
       \large Autoencoder Feature Extraction and XGBoost Classification\\
       on the UCI Taiwan Dataset}
\author{BigData 2026 Project Team}
\date{June 2026}

\begin{document}

\maketitle
\tableofcontents
\newpage

%─────────────────────────────────────────────────────────────────────────────
\section{Introduction}
%─────────────────────────────────────────────────────────────────────────────

The UCI Taiwan credit card dataset records payment history for 30,000 customers
from April to September 2005. The prediction target is whether a customer
defaulted on their October 2005 payment. Of the 30,000 customers, 6,636 (22.1\%)
defaulted, giving a 78/22 class split that makes the problem a standard
imbalanced binary classification task.

The project constraint fixes the model family: a Keras dense autoencoder
compresses 18 input features into a low-dimensional latent representation, and
XGBoost then classifies from those latents. We cannot switch to a different
classifier family, but we can control the autoencoder's training objective,
how latents are passed to XGBoost, how class imbalance is handled, and how
probabilities are post-processed into binary decisions.

This report covers every design decision in the pipeline: why each preprocessing
step exists, what the autoencoder and XGBoost do and why we configured them as
we did, what the evaluation metrics measure, and what seven ablation conditions
reveal about which parts of the architecture contribute to predictive quality.

%─────────────────────────────────────────────────────────────────────────────
\section{Dataset and Preprocessing}
%─────────────────────────────────────────────────────────────────────────────

\subsection{Raw features and hygiene}

The raw CSV contains 25 columns: a customer ID, the credit limit
(\texttt{LIMIT\_BAL}), three demographic variables (AGE, SEX, EDUCATION,
MARRIAGE), and 18 temporal columns recording, for each of the six months
April--September 2005, the repayment status (\texttt{PAY\_0}--\texttt{PAY\_6}),
the statement balance (\texttt{BILL\_AMT1}--\texttt{BILL\_AMT6}), and the
amount paid (\texttt{PAY\_AMT1}--\texttt{PAY\_AMT6}).

Three hygiene operations precede any feature engineering. First, the EDUCATION
column contains codes 0, 5, and 6 that the dataset README does not define; we
recode all three to 4, which the README labels ``others.'' Second, the MARRIAGE
column contains 54 rows with code 0, which is similarly undocumented; these are
recoded to 3 (``others''). Third, 35 rows are exact duplicates of other rows.
These are dropped before any train/test split, because leaving them in risks
placing the same record in both the training and test sets, inflating measured
performance. After deduplication, 29,965 rows remain.

\subsection{Feature engineering and column selection}

Feeding 18 raw temporal columns directly to an autoencoder has two costs:
high dimensionality and strong within-group multicollinearity. The six bill
amounts, for instance, track a single underlying credit exposure and correlate
with one another above $r = 0.9$. We engineer summary statistics that capture
the same information with fewer dimensions.

From the six payment-status columns we compute \texttt{Pay\_delay\_mean},
\texttt{Pay\_delay\_max}, and \texttt{Pay\_delay\_std}. From the six bill amounts
we compute \texttt{Bill\_mean}, \texttt{Bill\_std}, and \texttt{Bill\_max}. From
the six payment amounts we compute \texttt{Pays\_amts\_total} and
\texttt{Pays\_amts\_mean}. We also add two utilization features
(\texttt{Utilization\_1} and \texttt{Utilization\_mean}), a composite risk score,
and \texttt{n\_months\_over\_limit} (the count of months in which the outstanding
balance exceeded the credit limit). After engineering, all 18 raw temporal
columns except \texttt{PAY\_0} are dropped. \texttt{PAY\_0} records the most
recent payment status (September 2005) and is the single strongest predictor of
October default, with a Pearson correlation of 0.32 to the target.

Three engineered columns (\texttt{Pays\_amts\_total}, \texttt{Utilization\_1},
and \texttt{Risk\_score}) are kept in the saved CSV for interpretability by other
team members but are excluded from model training. \texttt{Pays\_amts\_total} is
exactly six times \texttt{Pays\_amts\_mean}; feeding both would introduce a
perfect linear dependency. \texttt{Utilization\_1} covers a single month and is
already summarised by \texttt{Utilization\_mean}. \texttt{Risk\_score} is a
linear combination of features already in the feature set. The final training
feature matrix has 18 columns.

\subsection{Log transformation of heavy-tailed monetary features}

The autoencoder is trained with mean squared error (MSE) loss. Before
StandardScaler normalisation, the monetary columns have right-skewed
distributions: credit limits span from 10,000 to 1,000,000 NT\$, and individual
bill amounts can reach 964,511 NT\$. After StandardScaler, one training row
reached $|z| = 72.9$ on a bill column. A single such outlier dominates the
gradient signal and forces the autoencoder to spend most of its capacity
reconstructing extreme rows at the cost of typical ones.

We apply a signed log-plus-one transform,
$\tilde{x} = \text{sign}(x) \cdot \log(1 + |x|)$,
to six columns: \texttt{LIMIT\_BAL}, \texttt{Bill\_mean}, \texttt{Bill\_std},
\texttt{Bill\_max}, \texttt{Pays\_amts\_mean}, and \texttt{Utilization\_mean}.
The signed form preserves the direction of negative values (negative bills arise
from refunds, which occur in about 2\% of rows). After transformation and
StandardScaler, the maximum $|z|$ across all features drops to the range 5--8.

\subsection{Categorical encoding}

SEX, EDUCATION, and MARRIAGE are nominal variables. The autoencoder reconstructs
each feature under MSE loss, which treats feature values as real numbers on a
continuous scale. For an ordinal-coded nominal, MSE imposes a false distance
structure: predicting MARRIAGE $= 2.3$ when the true code is 2 appears as a
small error even though the two values are semantically unrelated. We one-hot
encode all three columns with \texttt{drop\_first=True} to avoid perfect
multicollinearity, adding a net four binary columns (SEX: 1, EDUCATION: 2,
MARRIAGE: 1). \texttt{PAY\_0} is left numeric because it is genuinely ordinal:
higher values correspond to longer payment delays.

After all preprocessing steps the ML-ready file contains 29,965 rows and 22
columns (18 features plus the three interpretability-only columns and the target).

%─────────────────────────────────────────────────────────────────────────────
\section{Model Architecture}
%─────────────────────────────────────────────────────────────────────────────

\subsection{What an autoencoder does}

An autoencoder is a neural network trained to reproduce its own input at the
output layer. Internally it passes data through a bottleneck: a hidden layer
(the \emph{encoder}) with far fewer neurons than the input dimension. Because
the network must compress the 18-dimensional input through the bottleneck and
then reconstruct it, the bottleneck neurons are forced to capture the
information most necessary to represent the data. These bottleneck activations
are the \emph{latent representation}, and in this project they serve as the
features fed to XGBoost.

The motivation for this two-stage setup is that XGBoost receives a compact,
potentially de-noised representation of the input rather than raw correlated
features. In practice, as the ablation results show, the autoencoder does not
improve on feeding raw features directly to XGBoost on this dataset, a finding
we discuss in Section~\ref{sec:analysis}.

\subsection{Autoencoder architecture}

The autoencoder maps 18 input features through the sequence
$18 \to 24 \to 16 \to 10 \to 16 \to 24 \to 18$, where 10 is the bottleneck
dimension. The encoder path is:

\begin{center}
\begin{tabular}{lll}
\toprule
Layer & Units & Activation \\
\midrule
Input       & 18 & -- \\
Dense       & 24 & ReLU \\
Dropout     & -- & 0.2 drop rate \\
Dense       & 16 & ReLU \\
Dropout     & -- & 0.2 drop rate \\
Dense (bottleneck) & 10 & ReLU \\
\bottomrule
\end{tabular}
\end{center}

The decoder mirrors the encoder (Dense 16, Dense 24, Dense 18 with linear
activation). The output layer uses a linear activation because the targets are
continuous real values (StandardScaler outputs), not bounded to $[0, 1]$. The
network is trained with Adam and MSE loss, with early stopping on the validation
loss (patience = 5, restoring best weights).

\subsection{Weighted autoencoder loss (B.4)}

A class-blind MSE autoencoder reconstructs whatever the training distribution
presents. Because 78\% of training rows are non-defaulters, the autoencoder
implicitly allocates most of its capacity to reconstructing that majority. We
address this by passing \texttt{sample\_weight} to \texttt{autoencoder.fit}: each
defaulter row is up-weighted by $\frac{N_{\text{non-default}}}{N_{\text{default}}}
\approx 3.52$, the inverse class ratio. This forces the network to pay equal
attention to both classes in MSE terms, at no cost beyond one extra argument.
The ablation (Conditions C2 vs.\ C3) shows the effect is marginal: 0.0018
PR-AUC difference, within measurement noise.

\subsection{PAY\_0 concatenation (C.3)}

After encoding, we concatenate the scaled raw value of \texttt{PAY\_0} to the
10-dimensional latent vector, giving XGBoost an 11-dimensional input. The
rationale is that \texttt{PAY\_0} is the strongest predictor ($r = 0.32$), and
an autoencoder bottleneck of dimension 10 compressing 18 features may not
preserve this signal with full fidelity. Bypassing the bottleneck for this one
feature avoids that loss. The ablation confirms this matters: Condition C4
(latents only, no PAY\_0) achieves PR-AUC = 0.5123, the worst AE variant,
while C2 with PAY\_0 concatenated reaches 0.5217.

\subsection{XGBoost}

XGBoost is a gradient-boosted decision tree ensemble. It builds a sequence of
shallow trees, where each tree is trained to correct the residual errors of the
previous ensemble, and the final prediction is the sum of all trees' outputs
passed through a logistic function. ``Gradient boosting'' refers to fitting each
new tree to the gradient of the loss function with respect to the current
prediction.

For class imbalance, XGBoost provides \texttt{scale\_pos\_weight}: when set to
$w$, the positive-class gradient and hessian are multiplied by $w$, making each
positive example count as $w$ negative examples during tree construction. We
set $w = 3.52$.

We tune four hyperparameters via grid search: the number of trees
(200 or 400), tree depth (4 or 6), learning rate (0.05 or 0.10), and minimum
child weight (1 or 3). This gives 16 combinations. Each combination is evaluated
by 5-fold stratified cross-validation, scored on average precision (the area
under the precision-recall curve on the held-out fold). The grid consistently
selects shallow trees (max\_depth = 4), a conservative learning rate (0.05),
and the smaller tree count (200 estimators) across nearly all conditions and
seeds.

\subsection{Probability calibration (D.2)}

\texttt{scale\_pos\_weight} upweights the positive class during training, which
inflates the raw \texttt{predict\_proba} outputs: a score of 0.5 from a model
trained with $w = 3.52$ does not correspond to a 50\% probability of default.
Without correction, the Brier score (which depends on probability values, not
rankings) is meaningless.

We apply isotonic regression calibration using
\texttt{CalibratedClassifierCV(method='isotonic', cv='prefit')} on a 10\%
calibration slice carved from the training set before the autoencoder is
trained. Isotonic regression fits a non-decreasing step function mapping raw
model scores to empirical probabilities on the calibration data. The calibration
slice is held out from both the autoencoder's training and GridSearchCV's folds,
so the calibrator cannot overfit to data it has already influenced.

\subsection{Threshold selection (D.1)}

Binary predictions require a decision threshold: output 1 if
$P(\text{default}) \geq \tau$, else 0. The default $\tau = 0.5$ is wrong for
two reasons. First, after calibration the probabilities are honest but the
optimal operating point is problem-dependent, not always 0.5. Second, for a
22\%-positive dataset, the F1-maximising threshold is typically well below 0.5.

We select $\tau^*$ by computing the F1 score at every threshold on calibrated
calibration-slice predictions and taking the argmax:
\[
  \tau^* = \arg\max_\tau \, F1(\tau, \text{cal set})
\]
Using the calibration slice for this search prevents test-set leakage: $\tau^*$
is fixed before the test set is touched.

%─────────────────────────────────────────────────────────────────────────────
\section{Evaluation Metrics}
%─────────────────────────────────────────────────────────────────────────────

\subsection{Why ROC-AUC is not the primary metric}

The receiver operating characteristic (ROC) curve plots true positive rate
against false positive rate at every threshold. Its area (ROC-AUC) equals the
probability that a randomly drawn positive example is ranked above a randomly
drawn negative example. For a balanced dataset, ROC-AUC is informative. For an
imbalanced one, it can be misleading: a classifier that predicts ``no default''
for every row achieves a true positive rate of 0 and a false positive rate of 0,
which lies on the ROC diagonal and yields ROC-AUC close to 0.5. But it correctly
classifies all 78\% non-defaulters and fails on all 22\% defaulters. The metric
obscures this failure. We report ROC-AUC for reference only.

\subsection{Precision-Recall AUC (primary metric)}

The precision-recall (PR) curve plots precision ($\frac{TP}{TP+FP}$) against
recall ($\frac{TP}{TP+FN}$) at every threshold. PR-AUC is the area under this
curve. For an imbalanced problem, PR-AUC is more informative because precision
directly reflects how often a positive prediction is correct, regardless of how
many true negatives exist. A trivial all-negative classifier achieves
PR-AUC $\approx 0.22$ (the positive base rate) while a perfect classifier
achieves 1.0. All of our conditions score in the range 0.49--0.53, modest
improvement over the 0.22 baseline but well short of perfect discrimination.

\subsection{Brier score}

The Brier score is the mean squared error of predicted probabilities:
\[
  \text{Brier} = \frac{1}{N}\sum_{i=1}^N (p_i - y_i)^2
\]
where $p_i$ is the model's predicted probability and $y_i \in \{0, 1\}$ is the
true label. Lower is better. A model that always predicts the base rate
($p_i = 0.221$) achieves Brier $= 0.221 \times 0.779 = 0.172$. Our conditions
score between 0.135 and 0.140, beating this naive baseline. Brier score requires
well-calibrated probabilities; without isotonic calibration, the scores would be
inflated by \texttt{scale\_pos\_weight}.

\subsection{F1 at the optimal threshold}

F1 is the harmonic mean of precision and recall, $F1 = \frac{2 \cdot P \cdot R}{P + R}$.
Unlike PR-AUC, it depends on a fixed threshold. We report $F1@\tau^*$: the F1
score measured on the test set using the threshold $\tau^*$ selected on the
calibration slice. Because $\tau^*$ is tuned rather than fixed at 0.5, the score
reflects realistic operating performance rather than an arbitrary operating point.

\subsection{Accuracy}

Accuracy is the fraction of correct predictions. We report it for completeness,
but it is a poor summary of performance for a 22\%-positive dataset. A model
that predicts ``no default'' for every row achieves 77.9\% accuracy, which is
higher than some of our conditions' accuracy figures. The numbers in the results
table are accuracy at $\tau^*$, not at 0.5.

%─────────────────────────────────────────────────────────────────────────────
\section{Experiment Design}
%─────────────────────────────────────────────────────────────────────────────

We compare seven conditions, each run with three random seeds (seeds 0, 1, 2).
All conditions share the same 80/20 stratified train/test split (seed = 42), the
same 10\% calibration slice, the same 16-cell XGBoost grid, and the same
calibration and threshold-tuning protocol. Table~\ref{tab:conditions} defines
each condition and what it isolates.

\begin{table}[h]
\centering
\caption{Ablation conditions.}
\label{tab:conditions}
\begin{tabularx}{\textwidth}{llX}
\toprule
ID & Configuration & What it isolates \\
\midrule
C0 & XGB-only, $w = 3.52$ & XGB on raw 18 features, with imbalance correction \\
C1 & XGB-only, $w = 1$ & XGB on raw features, no imbalance correction \\
C2 & AE + XGB, weighted AE, PAY\_0 concat & Full current model (B.4 and C.3 active) \\
C3 & AE + XGB, class-blind AE, PAY\_0 concat & Effect of weighted AE loss (B.4) \\
C4 & AE + XGB, weighted AE, latents only & Effect of PAY\_0 concatenation (C.3) \\
C5 & AE + XGB, SMOTE on latents, $w = 1$ & Latent-space oversampling vs.\ $w$ \\
C6 & AE + XGB, Borderline-SMOTE on latents, $w = 1$ & Boundary-aware oversampling \\
\bottomrule
\end{tabularx}
\end{table}

For C0 and C1, XGBoost receives the 18 standardised raw features directly,
bypassing the autoencoder entirely. C0 uses \texttt{scale\_pos\_weight} = 3.52;
C1 sets it to 1 (no correction). These two conditions are the baseline against
which all AE variants are measured.

Conditions C2, C3, and C4 share the same autoencoder architecture (encoding\_dim
= 10) and differ only in training weight and feature concatenation:
\begin{itemize}
  \item C2 uses the weighted AE loss (B.4) and concatenates PAY\_0.
  \item C3 removes the weight (class-blind AE) but keeps PAY\_0 concatenation.
  \item C4 uses the weighted AE but does not concatenate PAY\_0.
\end{itemize}

C5 and C6 replace \texttt{scale\_pos\_weight} with synthetic oversampling applied
to the 10-dimensional latent space after encoding. C5 uses plain SMOTE, which
generates new minority samples by interpolating between existing ones in latent
space. C6 uses Borderline-SMOTE, which restricts synthetic sample generation to
latent vectors near the decision boundary. Both C5 and C6 use $w = 1$ to avoid
double-correcting for imbalance.

Three outer seeds per condition give 21 total runs. With only three seeds the
standard deviation estimate itself carries roughly 50\% relative error; a
difference smaller than one standard deviation cannot be treated as a reliable
ranking.

%─────────────────────────────────────────────────────────────────────────────
\section{Results}
%─────────────────────────────────────────────────────────────────────────────

Table~\ref{tab:results} reports the mean and standard deviation of each metric
across three seeds. The primary metric is PR-AUC.

\begin{table}[h]
\centering
\caption{Benchmark results: mean $\pm$ std across 3 seeds. Primary metric: PR-AUC.}
\label{tab:results}
\small
\begin{tabular}{lrrrrrrl}
\toprule
Cond. & PR-AUC & Brier & F1@$\tau^*$ & Accuracy & ROC-AUC & $\tau^*$ & Train (s) \\
\midrule
C0 & $0.5281 \pm 0.0065$ & $0.1350 \pm 0.0011$ & $0.5465 \pm 0.0054$ & $0.7959 \pm 0.0155$ & $0.7820 \pm 0.0077$ & 0.336 & 31 \\
C1 & $0.5280 \pm 0.0107$ & $0.1349 \pm 0.0011$ & $0.5463 \pm 0.0061$ & $0.7958 \pm 0.0150$ & $0.7831 \pm 0.0069$ & 0.343 & 31 \\
C2 & $0.5217 \pm 0.0046$ & $0.1362 \pm 0.0014$ & $0.5398 \pm 0.0045$ & $0.7680 \pm 0.0091$ & $0.7772 \pm 0.0089$ & 0.280 & 31 \\
C3 & $0.5235 \pm 0.0059$ & $0.1366 \pm 0.0013$ & $0.5367 \pm 0.0055$ & $0.7860 \pm 0.0103$ & $0.7764 \pm 0.0068$ & 0.317 & 31 \\
C4 & $0.5123 \pm 0.0027$ & $0.1382 \pm 0.0014$ & $0.5367 \pm 0.0066$ & $0.7825 \pm 0.0133$ & $0.7734 \pm 0.0075$ & 0.319 & 31 \\
C5 & $0.5177 \pm 0.0019$ & $0.1372 \pm 0.0009$ & $0.5389 \pm 0.0041$ & $0.7825 \pm 0.0070$ & $0.7743 \pm 0.0050$ & 0.316 & 44 \\
C6 & $0.4933 \pm 0.0029$ & $0.1398 \pm 0.0011$ & $0.5325 \pm 0.0051$ & $0.7581 \pm 0.0249$ & $0.7700 \pm 0.0052$ & 0.291 & 54 \\
\bottomrule
\end{tabular}
\end{table}

The top condition by PR-AUC is C0 at $0.5281 \pm 0.0065$. Four conditions are
within one standard deviation of C0: C1, C3, and C2. With only three seeds,
these four should be treated as statistically indistinguishable. The remaining
three conditions (C4, C5, C6) sit below one standard deviation from C0.

The $\tau^*$ column shows the F1-optimal threshold found on the calibration
slice. All conditions land in the range 0.28--0.34, well away from 0.5. The
isotonic calibration is functioning correctly: if calibration had failed and
raw scores were passed through, $\tau^*$ would be near 0.5 for the XGB-only
conditions and near 0 or 1 for the AE conditions. The AE conditions tend to
have lower $\tau^*$ (0.28--0.32) than the XGB-only conditions (0.34--0.34),
reflecting the different probability scales produced by the two model types
before calibration.

Training time is 31 seconds for all conditions that do not apply oversampling.
SMOTE (C5) adds 13 seconds and Borderline-SMOTE (C6) adds 23 seconds. The
autoencoder itself is not a bottleneck; at 21,572 training rows and 50 epochs
with early stopping, it fits in under 5 seconds on GPU.

%─────────────────────────────────────────────────────────────────────────────
\section{Analysis of Individual Conditions}
\label{sec:analysis}
%─────────────────────────────────────────────────────────────────────────────

\subsection{XGB-only beats every AE variant (C0, C1 vs.\ C2--C6)}

The best AE variant (C2: 0.5217) underperforms both XGB-only conditions (C0:
0.5281, C1: 0.5280) by 0.006 PR-AUC. This gap is small but consistent across
all three seeds. The autoencoder is adding information loss, not information
gain.

The likely cause is that the 18-feature dataset has high multicollinearity among
monetary columns, leaving little genuine non-linear structure for the autoencoder
to exploit. A direct empirical check supports this: a 2-neuron bottleneck trained
for 150 epochs achieves near-identical ROC-AUC to the 10-neuron bottleneck. When
a bottleneck that small can represent the data nearly as well as a 10-neuron one,
the autoencoder is performing PCA-equivalent linear compression, not capturing a
useful non-linear manifold. The model family is fixed by the project constraints;
this result is reported honestly.

\subsection{PAY\_0 concatenation is load-bearing (C4 vs.\ C2)}

Removing the PAY\_0 concatenation (C4: 0.5123) produces the worst AE variant,
0.0094 PR-AUC below C2 (0.5217). The AE bottleneck loses the strongest
predictor's signal, and without it XGBoost has to recover that information
indirectly from whatever the latent dimensions encode. It cannot fully compensate.
This is the most concrete design decision the ablation validates.

\subsection{Weighted AE loss has marginal effect (C2 vs.\ C3)}

C2 (weighted AE) scores 0.5217 vs.\ C3 (class-blind AE) at 0.5235. The
difference of 0.0018 is less than the standard deviation of either condition
(0.0046 and 0.0059, respectively) and is not distinguishable at three seeds.
Weighted AE loss is free to keep because it costs nothing at inference time and
has a principled motivation, but it is not the deciding factor.

\subsection{scale\_pos\_weight is not the lever for XGB-only (C0 vs.\ C1)}

C0 (with $w = 3.52$) scores 0.5281 and C1 (with $w = 1$) scores 0.5280 on
PR-AUC, a difference of 0.0001. The identical Brier scores (0.1350 vs.\ 0.1349)
confirm the two conditions are functionally the same. XGBoost does not require the imbalance correction on this dataset when
operating on raw features. The tree structure and PR-AUC scoring in
GridSearchCV already orient the model toward the minority class without
the explicit weight.

\subsection{Latent-space oversampling does not help (C5, C6)}

SMOTE on latents (C5: 0.5177) is slightly below the class-blind XGB baseline
C4 (0.5123 is C4; C5 improves on C4 slightly but both fall below C0/C1).
Borderline-SMOTE (C6: 0.4933) is the worst condition overall, 0.035 PR-AUC
below C0. Borderline-SMOTE restricts synthetic sample generation to latent
vectors near the classification boundary, which in a low-variance latent space
like this one means interpolating between points that are already near the
boundary on both sides, producing noisy synthetic samples that damage the
classifier's decision surface.

Plain SMOTE is neutral to mildly negative; Borderline-SMOTE is harmful.
Neither is recommended for this pipeline.

%─────────────────────────────────────────────────────────────────────────────
\section{Single-Run Kaggle Notebook Results}
%─────────────────────────────────────────────────────────────────────────────

The Kaggle notebook (\texttt{model.py}) was run at \texttt{encoding\_dim = 4}
(a deviation from the benchmark's \texttt{encoding\_dim = 10}). The results
from that run were:

\begin{center}
\begin{tabular}{ll}
\toprule
Metric & Value \\
\midrule
Best CV PR-AUC & 0.5425 \\
PR-AUC (test) & 0.5152 \\
Brier score & 0.1390 \\
F1 @ $\tau^* = 0.399$ & 0.5176 \\
Accuracy & 0.787 \\
ROC-AUC & 0.762 \\
\bottomrule
\end{tabular}
\end{center}

The test PR-AUC of 0.5152 is 0.0065 below the benchmark C2 result of 0.5217,
which uses \texttt{encoding\_dim = 10} on the same architecture. A 4-dimensional
bottleneck compresses 18 features into too few latent dimensions. The benchmark
confirms that 10 dimensions is better. The cross-validated training PR-AUC
(0.5425) is higher than the test PR-AUC (0.5152), a gap of 0.0273 that suggests
overfitting in the grid search with a smaller bottleneck.

%─────────────────────────────────────────────────────────────────────────────
\section{Limitations}
%─────────────────────────────────────────────────────────────────────────────

\paragraph{Three seeds are not enough to rank close conditions.}
Standard deviation estimates from three samples carry roughly 50\% relative
error. Four of the seven conditions lie within one standard deviation of the
top PR-AUC. No reliable ranking among those four is possible from this data.
Five or more seeds would be needed to confidently separate them.

\paragraph{The autoencoder is over-parameterised for this dataset.}
A 2-neuron bottleneck achieves near-identical AUC to a 10-neuron one. The
effective dimensionality of the data is low, consistent with the strong
multicollinearity in the monetary features. Using a deep autoencoder as a
feature extractor on data where PCA-equivalent linear projection captures
nearly all structure is adding architectural complexity without a corresponding
benefit.

\paragraph{The model family is fixed.}
XGBoost on raw features (C0) outperforms every AE+XGB variant. We cannot
act on this finding because the project constrains the classifier family. The
finding is real and is presented transparently.

\paragraph{Temporal data is treated as cross-sectional.}
Each row represents one customer's 6-month history. We do not model the time
series structure within or between rows. The dataset is a snapshot (April--September
2005) with a single binary outcome; extending the model to a rolling-window setup
would require a different experimental design.

%─────────────────────────────────────────────────────────────────────────────
\section{Recommendations}
%─────────────────────────────────────────────────────────────────────────────

\textbf{Revert \texttt{encoding\_dim} to 10 in the Kaggle notebook.}
The notebook was left at \texttt{encoding\_dim = 4} after an earlier
experiment. The benchmark shows dim = 10 achieves PR-AUC = 0.5217 vs.\ 0.5152
for dim = 4 on the same AE+XGB configuration. This is a one-line fix.

\textbf{Keep PAY\_0 concatenation (C.3).}
The ablation shows removing it drops PR-AUC by 0.0094. It is the most impactful
design choice in the AE+XGB pipeline.

\textbf{Keep isotonic calibration and threshold tuning.}
The calibrated threshold $\tau^*$ lands consistently near 0.30, not 0.5. Using
the default threshold would mismatch the model's probability scale and reduce
$F1@\tau^*$ by operating at the wrong point on the precision-recall curve.

\textbf{Do not use Borderline-SMOTE on latents.}
C6 is the worst overall condition, 0.035 PR-AUC below C0. Plain SMOTE (C5) is
neutral and adds 13 seconds of training time; it is not worth the complexity.

\textbf{If the model constraint were lifted,} XGB-only on the 18 standardised
features (C0) is the highest-performing setup in the benchmark. This should be
reported as the unconstrained recommendation.

\textbf{A one-class autoencoder approach (queued, not yet benchmarked)} trains
the AE exclusively on non-defaulters and uses per-sample reconstruction error as
an anomaly score feature. This is the theoretically correct use of an AE for
imbalanced classification and would be the next thing to test.

%─────────────────────────────────────────────────────────────────────────────
\section{Conclusion}
%─────────────────────────────────────────────────────────────────────────────

This project applies a Keras autoencoder and XGBoost pipeline to the UCI Taiwan
credit card default dataset, with a fixed model architecture constraint. The
preprocessing pipeline removes 35 duplicate rows, recodes undocumented category
codes in EDUCATION and MARRIAGE, engineers temporal summaries from six months of
payment data, and applies a signed log-plus-one transform to six heavy-tailed
monetary columns, bringing the maximum standardised z-score from 72.9 to
approximately 5--8.

The ablation benchmark (7 conditions, 3 seeds, 21 total runs) produces three
concrete findings. First, XGB-only on raw features outperforms every AE+XGB
variant by approximately 0.006 PR-AUC; the autoencoder is not contributing
positively, likely because the data's effective dimensionality is lower than
the bottleneck size. Second, concatenating raw \texttt{PAY\_0} with the latents
(C.3) is the most impactful single design choice within the AE+XGB family;
removing it costs 0.0094 PR-AUC. Third, Borderline-SMOTE on latents harms
performance, reducing PR-AUC by 0.035 relative to XGB-only.

The primary metric (PR-AUC 0.52--0.53 for the top conditions) shows modest
improvement over a naive majority-class baseline (PR-AUC $\approx 0.22$) and
meaningful improvement over random ranking (PR-AUC = 0.22), but clear room
remains for stronger feature engineering or, if the constraint were lifted, a
different model family.

\end{document}
